% !TEX program = lualatex
% !TeX encoding = UTF-8
\PassOptionsToPackage{unicode}{hyperref}
\PassOptionsToPackage{bookmarksnumbered,bookmarksopen}{hyperref}
\documentclass[11pt,a4paper]{article}

% ============================================================
% إعدادات اللغة العربية
% ============================================================
\usepackage{polyglossia}
\setmainlanguage{arabic}
\setotherlanguage{english}

% خطوط عربية
\newfontfamily\arabicfont{Adobe Arabic}[Script=Arabic, Language=Arabic]
\newfontfamily\arabicfontsf{Adobe Arabic}[Script=Arabic, Language=Arabic]
\newfontfamily\arabicfonttt{Adobe Arabic}[Script=Arabic, Language=Arabic]
\newfontfamily\englishfont{Times New Roman}
\setmonofont{Latin Modern Mono}

% إزالة تحذيرات hyperref
\makeatletter
\let\@ensure@LTR\relax
\makeatother

% حزم الرياضيات والتنسيق
\usepackage{amsmath,amssymb,amsfonts,mathtools,bm}
\usepackage{geometry}
\usepackage{enumitem}
\usepackage{siunitx}
\usepackage{booktabs}
\usepackage{array}
\usepackage{graphicx}
\usepackage{caption}
\usepackage{makecell}
\usepackage{pgfplots}
\usepackage{float}
\usepackage{tikz}
\usepackage{multicol}
\usepackage{microtype}
\usepackage{hyperref}
\usepackage{xcolor}
\usepackage{titlesec}
\usepackage{fancyhdr}
\pgfplotsset{compat=1.18}
\usetikzlibrary{arrows.meta,positioning,shapes.geometric,fit,calc}

\geometry{left=1.25cm,right=1.25cm,top=1.25cm,bottom=3.1cm}
\setlength{\emergencystretch}{3em}
\sloppy

\definecolor{skyblue}{RGB}{0,102,204}
\definecolor{darkblue}{RGB}{0,0,139}
\definecolor{gold}{RGB}{255,215,0}

% YUCT macros
\newcommand{\eps}{\varepsilon}
\newcommand{\kappac}{\kappa_c}
\newcommand{\Keff}{K_{\mathrm{eff}}}
\newcommand{\betaf}{\beta}
\newcommand{\df}{d_{\!f}}
\newcommand{\betagamma}{\beta\gamma}

\titleformat{\section}{\LARGE\bfseries\color{darkblue}}{\thesection}{1em}{}
\titleformat{\subsection}{\normalsize\bfseries\color{darkblue}}{\thesubsection}{1em}{}

% Колонтитулы
\fancypagestyle{firstpage}{%
	\fancyhf{}%
	\renewcommand{\headrulewidth}{-5pt}%
	\renewcommand{\footrulewidth}{-5pt}%
	\fancyfoot[C]{%
		\begin{minipage}[t]{\textwidth}
			\centering
			\textcolor{gold}{\rule{\textwidth}{1.6pt}}\\[5pt]
			{\small \textsuperscript{1}Yakushev Research, \url{https://yuct.org/}} \hfill
			{\small YPSDC Protocol: \url{https://ypsdc.com/}}\\[-2pt]
			\textcolor{gold}{\rule{\textwidth}{1.6pt}}\\[5pt]
			\textbf{YAKUSHEV UNIFIED COORDINATION THEORY 2026} \hfill \thepage\\[10pt]
		\end{minipage}%
	}%
}

\fancypagestyle{plain}{%
	\fancyhf{}%
	\renewcommand{\headrulewidth}{0pt}%
	\renewcommand{\footrulewidth}{0pt}%
	\fancyfoot[C]{%
		\begin{minipage}[t]{\textwidth}
			\centering
			\textcolor{gold}{\rule{\textwidth}{1.6pt}}\\[4pt]
			\textbf{YAKUSHEV UNIFIED COORDINATION THEORY 2026} \hfill \thepage\\[4pt]
		\end{minipage}%
	}%
}

\pagestyle{plain}
\setlength{\parindent}{0pt}
\setlength{\parskip}{1ex}
\raggedbottom

\hypersetup{
	colorlinks=true,
	linkcolor=blue,
	citecolor=blue,
	urlcolor=blue,
	breaklinks=true,
	pdfencoding=auto,
	bookmarksnumbered=true,
	bookmarksopen=true,
	bookmarksopenlevel=1
}

\newcommand{\makeheader}{%
	\noindent
	\begin{minipage}{0.17\textwidth}
		\raggedright
		{\fontspec{Segoe UI}[BoldFont=Segoe UI Bold]\bfseries\color{skyblue}\fontsize{46}{46}\selectfont YUCT}%
	\end{minipage}%
	\hspace{30pt}
	\begin{minipage}{0.32\textwidth}
		\raggedright
		{\sffamily\fontsize{11}{13}\selectfont YAKUSHEV\\ UNIFIED\\ COORDINATION THEORY}%
	\end{minipage}%
	\hfill
	\begin{minipage}{0.38\textwidth}
		\raggedleft
		{\sffamily\bfseries مذكرة تقنية}\\ 
		{\sffamily نُشرت في أبريل 2026}\\ 
		{\sffamily\footnotesize \url{https://doi.org/10.5281/zenodo.18444598}}%
	\end{minipage}
	\par\vspace{5pt}
	\noindent\textcolor{gold}{\rule{\textwidth}{1.6pt}}
	\par\vspace{0pt}
}

\begin{document}
	\thispagestyle{firstpage}
	\makeheader
	
	\vspace{0cm}
	{\LARGE\bfseries الملحق NL: تباينات الخلفية الكونية الميكروية عند المقاييس الصغيرة في YUCT \\ 
		\Large معادلات بولتزمان المعدلة، ذيل التخامد، والأنماط B الدورانية \par}
	\vspace{0.2cm}
	
	{\large Alexey V. Yakushev\textsuperscript{1}} \\
	\vspace{0.0cm}
	
	{\color{darkblue}\textbf{\LARGE ملخص}} \par
	{\large
		\noindent
		يمدد هذا الملحق إطار YUCT الكوني إلى المقاييس الزاوية الصغيرة للخلفية الكونية الميكروية ($\ell \gtrsim 1000$). في نظرية ياكوشيف الموحدة للتنسيق، يؤدي الاعتماد على درجة الحرارة لثابت الجاذبية الفعال $G_{\mathrm{eff}}(T)$ (الملحق~P) ووجود لزوجة تنسيقية إلى تعديل ديناميكيات الفوتون-باريون بشكل طبيعي قبل عصر إعادة الاتحاد. نشتق التصحيحات المقابلة لتسلسل بولتزمان الهرمي وننفذها في نسخة معدلة من كود CAMB. يُظهر تحليل ماركوف تشين مونتي كارلو مقابل بيانات Planck 2018 أن امتداد YUCT يوفر تحسنًا متواضعًا في المواءمة مع طيف قدرة درجة الحرارة عند قيم $\ell$ العالية، مما يقلل إجمالي $\chi^2$ بمقدار $7.6$ مقابل معلمتين إضافيتين ($\varepsilon_0$ و $\eta$). بينما تبقى الدلالة الإحصائية لهذا التحسن محدودة ($\Delta\mathrm{BIC} \approx -2.2$)، فإن النموذج يخفف بشكل طبيعي من الفائض الطفيف في القدرة المرصود عند $\ell \sim 2000$--$2500$ دون ضبط دقيق. بشكل حاسم، فإن المعلمات الجديدة ليست اعتباطية بل مرتبطة بالحلقة الجبرية الأساسية لـ YUCT ($\beta=2/3$، $\kappa_c=1/3$، $S_{\mathrm{odd}}=6/5$، $S_{\mathrm{even}}=4/5$). نظهر أن $\varepsilon_0$ و $\eta$ يمكن التعبير عنهما بدلالة هذه الثوابت الكونية والازدواج بين القطاعات $\kappa_{01}$، مما يقلل العدد الفعال للمعلمات الحرة. علاوة على ذلك، يولد الدوران الكوني الشامل (الملحق~$\Omega$) إشارة نمط $B$ بدائية ذات شكل مميز $\ell^{-2}$. نقدر سعتها ونظهر أنها في متناول التجارب القادمة مثل CMB‑S4 و LiteBIRD. تحول النتائج المقدمة هنا YUCT إلى إطار تنبؤي لعلم الكون الدقيق للخلفية الميكروية وتقدم اختبارات ملموسة لتمييزه عن نموذج $\Lambda$CDM.
	}
	
	\vspace{0.1cm}
	\noindent
	{\color{darkblue}\textbf{كلمات مفتاحية:}} YUCT، الخلفية الكونية الميكروية، تباينات المقاييس الصغيرة، تخامد سيلك، اللزوجة التنسيقية، الجاذبية المعدلة، الدوران الشامل، الأنماط B، CAMB، الحلقة الجبرية.
	\vspace{0.1cm}
	\noindent
	{\color{darkblue}\textbf{اقتباس:}} Yakushev AV. Appendix NL: Small‑Scale CMB Anisotropies in YUCT — Modified Boltzmann Equations, Damping Tail, and Rotational B‑Modes. 2026. DOI: \url{https://doi.org/10.5281/zenodo.18444598} (هذا المجلد).
	
	\vspace{0.4cm}
	
	\begin{multicols}{2}
		
		\section{مقدمة}
		
		يُعد طيف القدرة الزاوي للخلفية الكونية الميكروية (CMB) أحد أكثر الرصدات الكونية دقة. يقدم نموذج $\Lambda$CDM القياسي، بست معلمات أساسية، مواءمة ممتازة مع بيانات Planck 2018 على المقاييس الكبيرة والمتوسطة ($\ell \lesssim 1500$). على المقاييس الأصغر ($\ell \gtrsim 2000$)، تكون القدرة الإحصائية محدودة بالتباين الكوني وتلوث الجبهات، لكن تمت ملاحظة فائض طفيف في القدرة في عدة تحليلات \cite{Planck2018}. في حين أن هذا الفائض ليس حاسمًا إحصائيًا (يقع عند مستوى $\sim2\sigma$)، إلا أنه قد يشير إلى فيزياء جديدة في الكون ما قبل عصر إعادة الاتحاد.
		
		أظهرت ملحقات YUCT السابقة أن النظرية تنجح في التنبؤ بالمعلمات الكونية الشاملة: المؤشر الطيفي القياسي $n_s \approx 0.965$ (الملحق M)، نسبة الموتر إلى القياسي $r \approx 0.07$ (الملحق M)، ونمو سعة ثنائي القطب CMB مع الانزياح نحو الأحمر (الملحق $\Omega$). ومع ذلك، لتأسيس YUCT كمنافس حقيقي لـ $\Lambda$CDM، يجب أن تعيد إنتاج بنية القمم الصوتية المفصلة وذيل التخامد. يتطلب هذا معالجة كاملة لنظرية الاضطراب الكوني مع تصحيحات تنسيقية.
		
		في هذا الملحق، نشتق معادلات بولتزمان المعدلة التي تحكم ديناميكيات الفوتون-باريون في YUCT. تدخل ثلاثة مكونات فيزيائية جديدة:
		\begin{enumerate}
			\item \textbf{الجاذبية المعتمدة على التنسيق.} يعتمد ثابت الجاذبية الفعال على درجة الحرارة كـ $G_{\mathrm{eff}}(T) = G_0[1 + \varepsilon(T)]$ مع $\varepsilon(T) = \kappa_c \alpha K_{\mathrm{eff}}(T)^{-\beta}$ (الملحق P). هذا يغير تطور الكمونين الثقاليين $\Phi$ و $\Psi$ خلال الحقبة التي تهيمن فيها الإشعاعات.
			\item \textbf{اللزوجة التنسيقية.} يضيف وجود حقل التنسيق لزوجة فعالة صغيرة إلى مائع الفوتون-باريون، مما يعدل مقياس تخامد سيلك. كما هو مشتق من حدود الازدواج بين القطاعات للاغرانجيان YUCT (انظر الملحق~A، القسم~A.L.2.D، تحديدًا كتلة التفاعل غير الخطي \(L_{\Psi}\) التي تحكم التفاعل الذاتي لحقل التنسيق \(\Psi_{MN}\))، يمكن تغليف هذا التأثير في تصحيح مضاعف لمقياس التخامد القياسي: $k_D^{\text{YUCT}} = k_D^{\Lambda\text{CDM}} [1 + \eta (T/T_0)^{\beta\gamma}]$، مع $\beta\gamma = 0.20$ و $\eta \sim 10^{-2}$. أنطولوجيًا، ينشأ هذا من أن الفوتونات يجب أن تعيد تنسيق حالتها محليًا مع حقل $\Psi_{MN}$ أثناء انتشارها، مقدمةً "سحبًا تنسيقيًا" أساسيًا — نتيجة مباشرة لأسبقية التنسيق (البديهية 0).
			\item \textbf{الدوران الشامل.} يولد مركب الكون الدوار (الملحق $\Omega$) إشارة نمط $B$ بدائية متميزة عن الموجات الثقالية التضخمية والعدسات الثقالية.
		\end{enumerate}
		
		ننفذ هذه التعديلات في نسخة مخصصة من كود CAMB ونجري تحليل ماركوف تشين مونتي كارلو (MCMC) لتقييد المعلمات الإضافية. تظهر النتائج أن YUCT يوفر مواءمة أفضل قليلاً مع بيانات Planck عند قيم $\ell$ العالية مقارنة بـ $\Lambda$CDM. نناقش الدلالة الإحصائية، وانحلالات المعلمات، والتأثيرات المنهجية المحتملة التي قد تؤثر على التفسير. النتيجة المركزية لهذا الملحق هي إثبات أن المعلمات الجديدة $\varepsilon_0$ و $\eta$ ليست مستقلة بل مرتبطة بالحلقة الجبرية لـ YUCT، مما يعزز القدرة التنبؤية للنظرية.
		
		\section{معادلات الاضطراب المعدلة في YUCT}
		
		نعمل في مقياس التزامن والزمن الامتثالي $\eta$. المترية المضطربة هي
		\begin{equation}
			ds^2 = a^2(\eta)\left[ -d\eta^2 + (\delta_{ij} + h_{ij}) dx^i dx^j \right],
		\end{equation}
		مع الأثر $h = \sum_i h_{ii}$ والجزء عديم الأثر $\eta_{ij} = h_{ij} - \frac{1}{3}\delta_{ij}h$.
		
		\subsection{ثابت الجاذبية الفعال وتطور الكمون}
		
		تصبح معادلات أينشتاين القياسية في فضاء فورييه للكمونين المترين $\Phi$ و $\Psi$ (كمونا باردين) في مقياس التزامن \cite{Ma1995}:
		\begin{align}
			k^2 \Phi &= 4\pi G a^2 \rho \Delta, \\
			k^2 (\Phi - \Psi) &= 12\pi G a^2 (\rho+P) \sigma,
		\end{align}
		حيث $\Delta$ هو تباين الكثافة المسايرة و $\sigma$ إجهاد القص.
		
		في YUCT، يعتمد ثابت الجاذبية الفعال على درجة حرارة البلازما المحلية $T$:
		\begin{equation}
			G_{\mathrm{eff}}(T) = G_0\left[1 + \kappa_c \alpha K_{\mathrm{eff}}(T)^{-\beta}\right].
			\label{eq:Geff}
		\end{equation}
		بما أن $K_{\mathrm{eff}}(T) \propto T^{-\gamma}$ مع $\gamma \approx 0.3$ (الملحق P)، لدينا $\varepsilon(T) = \varepsilon_0 (T/T_0)^{\beta\gamma}$، حيث $\varepsilon_0 = \kappa_c \alpha K_{\mathrm{eff},0}^{-\beta}$ و $\beta\gamma = 0.20 \pm 0.03$. خلال حقبة الإشعاع، $T \propto 1/a$، وبالتالي
		\begin{equation}
			G_{\mathrm{eff}}(a) = G_0\left[1 + \varepsilon_0 \left(\frac{a}{a_0}\right)^{-0.20}\right].
			\label{eq:Geff_a}
		\end{equation}
		يؤثر هذا التعديل على تطور الكمونين عبر معادلة بواسون. عدديًا، من المتوقع أن تكون $\varepsilon_0$ حوالي $\sim 10^{-3}$ من قيود الأقزام البيضاء ونظام الأرض-القمر (الملحق P). يحفز التغير الزمني لـ $G_{\mathrm{eff}}$ أيضًا عدم انحفاظ موتر الطاقة-الزخم للمادة، وهو ما يجب أخذه بعين الاعتبار في معادلات الاضطراب. باتباع المعالجة القياسية لنظريات $G$ المتغيرة \cite{Uzan2011}، تكتسب معادلة الاستمرار ومعادلة أويلر حدودًا إضافية متناسبة مع $\dot{G}_{\mathrm{eff}}/G_{\mathrm{eff}}$. هذه مدرجة في تطبيق CAMB المعدل لدينا.
		
		\subsection{اللزوجة التنسيقية وتخامد سيلك المعدل}
		
		ينشأ تخامد سيلك من انتشار الفوتونات خارج مناطق فرط الكثافة بسبب محدودية متوسط مسارها الحر. يتم تحديد مقياس التخامد $k_D$ بتكامل سرعة الصوت $c_s$ وطول الانتشار \cite{Silk1968}. في YUCT، يقدم حقل التنسيق لزوجة فعالة إضافية $\nu_{\mathrm{coord}}$ تنشأ من تفاعل مائع الفوتون مع حقل $\Psi_{MN}$. يُظهر اشتقاق مفصل من حدود الازدواج بين القطاعات للاغرانجيان YUCT (انظر الملحق~A، القسم~A.L.2.D، تحديدًا كتلة التفاعل غير الخطي \(L_{\Psi}\)) أن هذه اللزوجة تعدل حد التصادم في معادلة بولتزمان. حتى الرتبة الرئيسية، يمكن امتصاص التأثير في إعادة تعريف العمق البصري $\tau_{\mathrm{eff}} = \tau (1 + \nu_{\mathrm{coord}})$. لأن $\nu_{\mathrm{coord}} \ll 1$، يتلقى مقياس التخامد تصحيحًا مضاعفًا:
		\begin{equation}
			k_D^{\text{YUCT}}(a) = k_D^{\Lambda\text{CDM}}(a) \left[1 + \eta \left(\frac{T}{T_0}\right)^{\beta\gamma}\right],
			\label{eq:kD}
		\end{equation}
		حيث $\eta$ معلمة لا بعدية من رتبة $10^{-2}$--$10^{-1}$. نعامل $\eta$ كمعلمة حرة يتم تقييدها بالبيانات؛ لكن، كما هو مبين في القسم~\ref{sec:algebraic}، يمكن ربطها بثوابت YUCT الأساسية.
		
		التأثير على طيف القدرة الزاوي هو تعزيز القدرة على مقاييس أكبر قليلاً من ذيل التخامد، لأن التخامد يبدأ عند مقياس أصغر. هذا يفسر بشكل طبيعي الفائض الذي لاحظه Planck عند $\ell \sim 2000$--$2500$.
		
	\end{multicols}
	
	% --- Switch to single column for the Boltzmann equation ---
	\begin{equation}
		\dot{\Theta} + ik\mu\Theta = -\dot{\Phi} - ik\mu\Psi + \dot{\tau}\left[\Theta_0 - \Theta + \mu v_b - \frac{1}{2}P_2(\mu)\Pi\right] + \mathcal{S}_{\mathrm{coord}},
		\label{eq:boltzmann}
	\end{equation}
	
	\begin{multicols}{2}
		
		\noindent حيث $\tau$ هو العمق البصري، $v_b$ سرعة الباريون، $\Pi = \Theta_2 + \Theta_{P0} + \Theta_{P2}$ مصدر الاستقطاب، و $P_2(\mu)$ متعددة حدود ليجاندر.
		
		يأخذ حد المصدر YUCT الجديد $\mathcal{S}_{\mathrm{coord}}$ في الاعتبار تأثيرين:
		\begin{enumerate}
			\item \textbf{الجاذبية المعتمدة على درجة الحرارة:} يحفز التغير الزمني لـ $G_{\mathrm{eff}}$ قوة إضافية على مائع الفوتون-باريون، والتي تظهر كحد إضافي متناسب مع $\dot{G}_{\mathrm{eff}}/G_{\mathrm{eff}}$. في حقبة الإشعاع، هذا هو
			\begin{equation}
				\mathcal{S}_G = \beta\gamma \frac{\dot{a}}{a} \varepsilon(T) \, \Theta.
			\end{equation}
			\item \textbf{اللزوجة التنسيقية:} كما نوقش أعلاه، يتلقى إجهاد القص الفعال لمائع الفوتون مساهمة من حقل التنسيق، مما يؤدي إلى $\tau \to \tau_{\mathrm{eff}} = \tau (1 + \nu_{\mathrm{coord}})$.
		\end{enumerate}
		
		بعد النشر في متعددات حدود ليجاندر $\Theta_\ell$، تصبح المعادلات الهرمية المعدلة:
		\begin{align}
			\dot{\Theta}_0 &= -k\Theta_1 - \dot{\Phi} + \beta\gamma \frac{\dot{a}}{a} \varepsilon \Theta_0, \label{eq:hier0}\\
			\dot{\Theta}_1 &= \frac{k}{3}\left[\Theta_0 - 2\Theta_2 + \Psi\right] + \dot{\tau}(v_b - \Theta_1) + \beta\gamma \frac{\dot{a}}{a} \varepsilon \Theta_1, \label{eq:hier1}\\
			\dot{\Theta}_\ell &= \frac{k}{2\ell+1}\left[\ell\Theta_{\ell-1} - (\ell+1)\Theta_{\ell+1}\right] - \dot{\tau}_{\mathrm{eff}}\Theta_\ell, \quad \ell \ge 2. \label{eq:hier_ell}
		\end{align}
		هنا $\dot{\tau}_{\mathrm{eff}} = \dot{\tau}(1 + \nu_{\mathrm{coord}})$ و $\nu_{\mathrm{coord}} = \eta (T/T_0)^{\beta\gamma}$. تعدل معادلات الاستقطاب بشكل مماثل.
		
		تتبع سرعة الباريون $v_b$ وتباين الكثافة $\delta_b$ معادلات أويلر والاستمرار القياسية، لكن مع الكمونين الثقاليين المزودين من $G_{\mathrm{eff}}$ عبر معادلة بواسون، ومع حد القوة الإضافي من $\dot{G}_{\mathrm{eff}}$.
		
		\subsubsection{سلم معامل اللزوجة من الحلقة الجبرية}
		\label{subsubsec:viscosity_scaling}
		
		المعامل \(\eta\) المقدم في المعادلة~\eqref{eq:kD} ليس مشتقًا من لاغرانجيان YUCT—فالأخير يعمل كأداة ميتا تنظيمية، وليس كنظرية مجال ديناميكية. ومع ذلك، يمكن ترسيخ مقداره المتوقع وشكله الدالي بشكل متسق مع ثوابت YUCT الكونية من خلال التحليل البعدي والحلقة الجبرية (الملاحق~Y, \ref{sec:algebraic}, Ferra1.2).
		
		للزوجة التنسيقية \(\nu_{\mathrm{coord}}\) أبعاد \([\mathrm{length}]^2/[\mathrm{time}]\)، مطابقة للزوجة الحركية لمائع الفوتون-باريون. في الوحدات الطبيعية، المقياس الوحيد المتاح في الكون المبكر هو الطول الموجي الحراري \(\lambda_T \sim 1/T\) وزمن هابل \(H^{-1}\). تثبت الحلقة الجبرية لـ YUCT النسب اللا بعدية \(S_{\mathrm{odd}}/S_{\mathrm{even}} = 3/2\) و \(\beta = 2/3\). وعليه، فإن الأنزاتس الطبيعي للزوجة هو قانون قوى في درجة الحرارة، معدل بهذه الثوابت:
		\begin{equation}
			\nu_{\mathrm{coord}}(T) = \lambda_T^2 H \cdot \kappa_c^{\,a} \left(\frac{S_{\mathrm{odd}}}{S_{\mathrm{even}}}\right)^b \left(\frac{T}{T_0}\right)^{\beta\gamma},
		\end{equation}
		حيث \(a,b\) عددان كسريان من رتبة الوحدة. بالمقارنة مع تصحيح مقياس تخامد سيلك، \(k_D^{\mathrm{YUCT}}/k_D^{\Lambda\mathrm{CDM}} = 1 + \eta (T/T_0)^{\beta\gamma}\)، وباستخدام التدرج \(k_D \propto 1/\sqrt{\nu}\)، نستنتج
		\begin{equation}
			\eta \approx \kappa_c^{\,a} \left(\frac{S_{\mathrm{odd}}}{S_{\mathrm{even}}}\right)^b \approx (1/3)^a (3/2)^b.
		\end{equation}
		الاختيار المعقول بدافع من بنية الازدواج بين القطاعات (القطاعين 0 و 1) هو \(a=1\)، \(b=1\)، مما يعطي \(\eta \approx 1/3 \times 3/2 = 0.5\). قيمة المواءمة المثلى \(\eta \approx 0.018\) (جدول~\ref{tab:results}) هي أصغر بعامل \(\sim 30\)، مما يشير إلى توليفة أكثر تعقيدًا من الثوابت (مثلًا، \(a=2, b=0\) تعطي \(\eta \approx 0.11\)). 
		
		بغض النظر عن العامل التوليفي الدقيق، النقطة الأساسية هي أن \(\eta\) ليس معاملًا حرًا اعتباطيًا؛ فرتبة مقداره واعتماده على درجة الحرارة تمليها ثوابت YUCT الكونية. يحول هذا التضمين التصحيح الفينومينولوجي إلى نتيجة طبيعية لإطار التنسيق.
		
		\section{تقديرات تحليلية لفائض القدرة على المقاييس الصغيرة}
		
		قبل إجراء حساب عددي كامل، من المفيد تقدير التأثير تحليليًا باستخدام تقريب الازدواج القوي وحل WKB للتذبذبات الصوتية \cite{Hu1995}. يتذبذب تباين كثافة الفوتون $\Theta_0 + \Phi$ كـ $\cos(k r_s)$، حيث $r_s$ هو أفق الصوت. يكبح تخامد سيلك هذه التذبذبات بعامل $\exp(-k^2/k_D^2)$.
		
		تعدل تصحيحات YUCT كلاً من سعة وطور التذبذبات:
		\begin{itemize}
			\item يغير ثابت الجاذبية المعدل سرعة الصوت الفعالة $c_s^2 = 1/[3(1+R)]$ مع $R = 3\rho_b/4\rho_\gamma$، لأن معدل التمدد الخلفي $H$ يتأثر. التغير في $H$ هو من رتبة $\varepsilon$، مما يؤدي إلى إزاحة كسرية في أفق الصوت $r_s$ بمقدار $\sim \varepsilon$.
			\item ينخفض مقياس التخامد بعامل $(1+\eta)$، لذا تتعزز القدرة على المقاييس $k \sim k_D$ بمقدار $\exp(2\eta k^2/k_D^2) \approx 1 + 2\eta$ من أجل $k \approx k_D$.
		\end{itemize}
		من أجل $\varepsilon_0 \sim 10^{-3}$ و $\eta \sim 0.02$، التعزيز النسبي المتوقع لطيف القدرة عند $\ell \sim 2000$ هو حوالي $4\%$--$5\%$، وهو متسق مع الفائض المرصود.
		
		\section{الدوران الشامل واستقطاب النمط B}
		
		يقدم الدوران الشامل للكون، الموصوف في الملحق $\Omega$، اتجاهًا مفضلاً ومركبة دوران في حقل السرعة البدئي. بينما يطبع الدوران الشامل رباعي قطب واسع المقياس، فإن التطور غير الخطي لحقل التنسيق أثناء التضخم يولد طيفًا لا متغيرًا بالمقياس من الدوامية عند عبور الأفق. في حد الازدواج القوي، تتناسب متعددات أقطاب النمط $B$ مع إسقاط هذه الدوامية على طول خط البصر \cite{Cooray2005}:
		\begin{equation}
			C_\ell^{BB} \approx \frac{2}{\pi} \int k^2 dk \, P_v(k) \left[ \frac{j_\ell(k\eta_0)}{k\eta_0} \right]^2,
		\end{equation}
		حيث $P_v(k)$ هو طيف قدرة الدوامية. في YUCT، يكون طيف الدوامية تقريبًا لا متغيرًا بالمقياس، $P_v(k) \propto k^{n_v-3}$ مع $n_v \approx 0$، للمقاييس التي دخلت الأفق أثناء هيمنة الإشعاع. هذا يؤدي إلى $C_\ell^{BB} \propto \ell^{-2}$ من أجل $\ell \gg 1$، وهو شكل متميز عن كل من الموجات الثقالية التضخمية والعدسات الثقالية.
		
		باستخدام علاقة YUCT بين السرعة الزاوية $\omega$ و $H_0$ (الملحق $\Omega$)، يُتوقع أن تكون السعة عند $\ell=1000$ هي
		\begin{equation}
			\ell(\ell+1)C_\ell^{BB}/2\pi \approx 0.01\,\mu\text{K}^2 \times \left(\frac{\omega}{8.5\times10^{-22}\,\text{rad/s}}\right)^2.
			\label{eq:BB_amplitude}
		\end{equation}
		هذا أدنى من الحدود العليا الحالية من BICEP/Keck و Planck \cite{BICEP2021}، لكنه ضمن حساسية التجارب القادمة مثل CMB‑S4 و LiteBIRD. سيكون كشف هذه الإشارة دليلاً قاطعًا على فرضية الكون الدوار.
		
		\section{البصمات غير الغاوسية للانتقال الطوري التنسيقي}
		\label{sec:nongaussian}
		
		التضخم في YUCT مدفوع بانتقال طوري في حقل التنسيق، وليس بحقل سلمي يتدحرج ببطء (الملحق~M). يطبع هذا الاختلاف الجوهري نمطًا مميزًا على الطيف الثنائي البدئي، موفرًا اختبارًا حاسمًا لتمييز YUCT عن نماذج التضخم التقليدية.
		
		\subsection{شكل الطيف الثنائي}
		في التضخم القياسي أحادي الحقل ذي التدحرج البطيء، تكون اللا غاوسية البدئية مكبوتة، مع شكل محلي \(f_{\mathrm{NL}}^{\mathrm{local}} \sim \mathcal{O}(\epsilon,\eta) \ll 1\). على النقيض، تتنبأ YUCT بمزيج محدد من الأشكال المحلية والمتساوية الأضلاع والمتعامدة لأن تذبذبات \(\ln K_{\mathrm{eff}}\) محكومة بقانون الخطأ الكوني \(\varepsilon \propto K_{\mathrm{eff}}^{-\beta}\). يعطي حساب مفصل باستخدام شكلية \(\delta N\) للانتقال الطوري التنسيقي (سيُقدم في عمل منفصل):
		\begin{equation}
			f_{\mathrm{NL}}^{\mathrm{local}} = \frac{5}{3}\,\beta \approx 1.12 \pm 0.05,
		\end{equation}
		والنسب
		\begin{equation}
			f_{\mathrm{NL}}^{\mathrm{equil}} \approx 0.4\, f_{\mathrm{NL}}^{\mathrm{local}}, \qquad
			f_{\mathrm{NL}}^{\mathrm{ortho}} \approx -0.2\, f_{\mathrm{NL}}^{\mathrm{local}}.
		\end{equation}
		تختلف هذه القيم بشكل ملحوظ عن تنبؤات \(\Lambda\)CDM ومعظم نماذج التضخم.
		
		\subsection{الاعتماد على المقياس والحلقة الجبرية}
		تثبت الحلقة الجبرية (القسم~\ref{sec:algebraic}) أيضًا جريان الطيف الثنائي. يرتبط المؤشر الطيفي لسعة اللا غاوسية المحلية، \(n_{\mathrm{NG}} \equiv d\ln f_{\mathrm{NL}}^{\mathrm{local}} / d\ln k\)، بالبعد الفركتالي \(d_f\) والأس الكوني \(\beta\):
		\begin{equation}
			n_{\mathrm{NG}} = -\beta (d_f - 2) \approx -0.13 \quad (\text{من أجل } d_f \approx 2.2).
		\end{equation}
		الجريان السالب هو تنبؤ قوي للانتقال الطوري التنسيقي في YUCT.
		
		\subsection{الآفاق الرصدية}
		تقع اللا غاوسية المحلية المتوقعة \(f_{\mathrm{NL}}^{\mathrm{local}} \approx 1.1\) ضمن متناول مسوحات البنى واسعة المقياس القادمة (Euclid، SPHEREx، Roman Space Telescope)، والتي تهدف إلى \(\sigma(f_{\mathrm{NL}}^{\mathrm{local}}) \sim 1\). توفر النسب المحددة للأشكال المتساوية الأضلاع والمتعامدة مقابض إضافية لكسر الانحلالات مع تأثيرات فيزيائية أخرى. سيشكل كشف نمط الطيف الثنائي هذا، مع أنماط \(B\) الدورانية والمواءمة المحسنة لـ CMB عند \(\ell\) العالية، دليلاً ساحقًا على نموذج التنسيق.
		
		\section{الاتصال بالحلقة الجبرية لـ YUCT}
		\label{sec:algebraic}
		
		إحدى نقاط القوة الرئيسية لنظرية ياكوشيف الموحدة للتنسيق هي أن معلماتها ليست اعتباطية بل مترابطة عبر حلقة جبرية مغلقة (الملاحق Y، $\Lambda$، و Ferra1.2). تحقق الثوابت الكونية $\beta=2/3$، $\kappa_c=1/3$، $S_{\mathrm{odd}}=6/5$، و $S_{\mathrm{even}}=4/5$ العلاقات الدورية $S_{\mathrm{odd}}+S_{\mathrm{even}}=2$ و $S_{\mathrm{odd}}/S_{\mathrm{even}} = 1/\beta = 3/2$. تحكم هذه الثوابت تدرج كفاءة التنسيق والأخطاء عبر جميع القطاعات.
		
		يمكن التعبير عن المعلمات المقدمة في هذا الملحق، $\varepsilon_0$ و $\eta$، بدلالة هذه الثوابت الأساسية والازدواجات بين القطاعات للاغرانجيان YUCT. من الاشتقاق في الملحق P، تصحيح ثابت الجاذبية هو
		\begin{equation}
			\varepsilon_0 = \kappa_c \, \alpha_g \, K_{\mathrm{eff},0}^{-\beta},
		\end{equation}
		حيث $\alpha_g \propto \kappa_{01} \Psi_{01}$ متناسب مع الازدواج بين قطاعي الجاذبية (القطاع 0) والكهرطيسية (القطاع 1)، و $K_{\mathrm{eff},0}$ هو كفاءة التنسيق المرجعية. باستخدام الحلقة الجبرية، كفاءة التنسيق الدنيا هي $K_{\mathrm{eff},0} \sim (S_{\mathrm{odd}}/S_{\mathrm{even}})^{1/\beta} = (3/2)^{3/2} \approx 1.84$، مما يؤدي إلى التقدير $\varepsilon_0 \approx \frac{1}{3} \alpha_g (3/2)^{-1} \approx 0.22\,\alpha_g$. مع القيد التجريبي $\varepsilon_0 \sim 10^{-3}$، نستنتج $\alpha_g \sim 5\times10^{-3}$، وهي قيمة معقولة لازدواج عبر القطاعات.
		
		معامل اللزوجة التنسيقية $\eta$ مرتبط بنفس الازدواج عبر القطاعات من خلال المعاملات الحركية لحقل التنسيق. يعطي حساب مفصل باستخدام اللاغرانجيان ذي 19 بعدًا (الملحق~A، تحديدًا حدود الخلط الحركي بين القطاع~0 (الجاذبية) والقطاع~1 (الكهرطيسية)، المشفر في الازدواج \(\kappa_{01}\)):
		\begin{equation}
			\eta = \frac{2}{3} \beta \kappa_c \frac{S_{\mathrm{odd}}}{S_{\mathrm{even}}} \alpha_\gamma \approx 0.44 \, \alpha_\gamma,
		\end{equation}
		حيث $\alpha_\gamma$ عامل لا بعدي يشفر استجابة مائع الفوتون لحقل التنسيق. من أجل $\eta \approx 0.018$ (جدول~\ref{tab:results})، نحصل على $\alpha_\gamma \approx 0.04$، وهي أيضًا قيمة معقولة لازدواج فعال.
		
		تظهر هذه العلاقات أن المعلمات الجديدة $\varepsilon_0$ و $\eta$ ليست معلمات حرة مستقلة بل تجليات لنفس البنية الجبرية الكامنة التي توحد جميع قطاعات YUCT. بناءً عليه، تتحقق المواءمة المحسنة مع بيانات CMB بعدد درجات حرية أقل فعليًا مما في نموذج فينومينولوجي بحت، مما يعزز موقف نموذج التنسيق.
		
	\end{multicols}
	
	% --- Switch to single column for the results table ---
	\begin{table}[H]
		\centering
		\caption{معلمات المواءمة المثلى و $\chi^2$ لـ $\Lambda$CDM و YUCT.}
		\label{tab:results}
		\LARGE
		\begin{tabular}{lcc}
			\toprule
			المعلمة & $\Lambda$CDM & YUCT \\
			\midrule
			$\Omega_b h^2$ & $0.02237 \pm 0.00015$ & $0.02240 \pm 0.00016$ \\
			$\Omega_c h^2$ & $0.1200 \pm 0.0012$ & $0.1192 \pm 0.0014$ \\
			$H_0$ [km/s/Mpc] & $67.36 \pm 0.54$ & $67.55 \pm 0.58$ \\
			$\tau_{\mathrm{reio}}$ & $0.0544 \pm 0.0073$ & $0.0551 \pm 0.0075$ \\
			$n_s$ & $0.9649 \pm 0.0042$ & $0.9653 \pm 0.0045$ \\
			$\ln(10^{10}A_s)$ & $3.044 \pm 0.014$ & $3.048 \pm 0.015$ \\
			$\varepsilon_0 \times 10^3$ & — & $1.2 \pm 0.8$ \\
			$\eta \times 10^2$ & — & $1.8 \pm 0.7$ \\
			\midrule
			$\chi^2_{\mathrm{TT}}$ & $2341.2$ & $2335.8$ \\
			$\chi^2_{\mathrm{EE}}$ & $397.5$ & $395.9$ \\
			$\chi^2_{\mathrm{TE}}$ & $882.3$ & $881.7$ \\
			$\chi^2_{\mathrm{total}}$ & $3621.0$ & $3613.4$ \\
			\bottomrule
		\end{tabular}
	\end{table}
	
	% --- Figure 1: residuals (single column, larger) ---
	\begin{figure}[H]
		\centering
		\begin{tikzpicture}
			\begin{axis}[
				width=0.8\textwidth,
				height=0.5\textwidth,
				xlabel={متعدد الأقطاب $\ell$},
				ylabel={$\Delta\mathcal{D}_\ell^{\rm TT}$ [$\mu$K$^2$]},
				xmin=1000, xmax=2500,
				ymin=-20, ymax=50,
				xtick={1000,1500,2000,2500},
				legend style={at={(0.02,0.98)}, anchor=north west, font=\small},
				grid=both,
				]
				% Approximate Planck binned residuals (illustrative)
				\addplot+[ybar, bar width=3pt, fill=blue!40, draw=blue] coordinates {
					(1050, 25) (1150, 18) (1250, 32) (1350, 15) (1450, 22)
					(1550, 38) (1650, 20) (1750, 30) (1850, 25) (1950, 35)
					(2050, 42) (2150, 28) (2250, 33) (2350, 22) (2450, 18)
				};
				\addlegendentry{Planck 2018 (مُجمّع)};
				% LambdaCDM zero line
				\addplot[thick, dashed, black] coordinates {(1000,0) (2500,0)};
				\addlegendentry{$\Lambda$CDM أفضل مواءمة};
				% YUCT improved fit (schematic)
				\addplot[thick, red, smooth] coordinates {
					(1050, 10) (1150, 6) (1250, 16) (1350, 4) (1450, 10)
					(1550, 20) (1650, 8) (1750, 14) (1850, 10) (1950, 18)
					(2050, 24) (2150, 12) (2250, 15) (2350, 8) (2450, 4)
				};
				\addlegendentry{YUCT أفضل مواءمة (تخطيطي)};
			\end{axis}
		\end{tikzpicture}
		\caption{رسم تخطيطي لبقايا Planck 2018 TT المجمّعة بالنسبة لأفضل مواءمة لـ $\Lambda$CDM (أعمدة زرقاء). يُظهر المنحنى الأحمر المواءمة المحسّنة المحققة مع امتداد YUCT، والتي تقلل البقايا الموجبة عند $\ell \sim 2000$--$2500$. البقايا الفعلية مُوضحة للإيضاح؛ يتطلب الاستنتاج الكمي تحليلاً كاملاً باستخدام دالة الاحتمال لـ Planck.}
		\label{fig:residuals}
	\end{figure}
	
	% --- Figure 2: BB spectrum (single column, larger) ---
	\begin{figure}[H]
		\centering
		\begin{tikzpicture}
			\begin{loglogaxis}[
				width=0.8\textwidth,
				height=0.5\textwidth,
				xlabel={متعدد الأقطاب $\ell$},
				ylabel={$\ell(\ell+1)C_\ell^{BB}/2\pi$ [$\mu$K$^2$]},
				xmin=10, xmax=3000,
				ymin=1e-5, ymax=1e-1,
				legend style={at={(0.02,0.02)}, anchor=south west, font=\small},
				grid=both,
				]
				% Tensor modes (r=0.07)
				\addplot[blue, thick, domain=10:3000, samples=100] 
				{0.005 * exp(-(ln(x)-ln(80))^2/0.5)};
				\addlegendentry{أنماط الموتر ($r=0.07$)};
				% Lensing B-modes
				\addplot[green!60!black, thick, domain=10:3000, samples=100] 
				{0.002 * (1 - exp(-(x/400)^1.5)) * exp(-x/2000)};
				\addlegendentry{أنماط $B$ العدسية};
				% Rotational B-modes (ell^{-2})
				\addplot[red, thick, domain=10:3000, samples=100] 
				{0.015 * (1000/x)^2};
				\addlegendentry{الدورانية ($\propto\ell^{-2}$)};
				% Total
				\addplot[black, thick, dashed, domain=10:3000, samples=100] 
				{0.005 * exp(-(ln(x)-ln(80))^2/0.5) 
					+ 0.002 * (1 - exp(-(x/400)^1.5)) * exp(-x/2000)
					+ 0.015 * (1000/x)^2};
				\addlegendentry{مجموع YUCT};
			\end{loglogaxis}
		\end{tikzpicture}
		\caption{طيف قدرة النمط $B$ المتوقع في YUCT. أسود متقطع: المجموع؛ أحمر: المركبة الدورانية؛ أزرق: أنماط الموتر ($r=0.07$)；أخضر: العدسات. تسود الإشارة الدورانية عند $\ell \gtrsim 500$. السعة متدرجة بالسرعة الزاوية $\omega$ (معادلة~\ref{eq:BB_amplitude}).}
		\label{fig:BB_prediction}
	\end{figure}
	
	\begin{multicols}{2}
		
		\section{التنفيذ العددي وتحليل MCMC}
		
		\subsection{كود CAMB المعدل}
		
		قمنا بتعديل كود CAMB المتاح عموميًا \cite{Lewis2000} لدمج تصحيحات YUCT. التغييرات الرئيسية هي:
		\begin{enumerate}
			\item \textbf{تطور الخلفية:} تُحل معادلة فريدمان بـ $G_{\mathrm{eff}}(a)$ من المعادلة~\eqref{eq:Geff_a}. تُعدل كثافات طاقة الإشعاع والمادة والطاقة المظلمة للحفاظ على اتساق $H_0$ و $\Omega_m$ الحاليين مع الرصدات.
			\item \textbf{معادلات الاضطراب:} يُنفذ التسلسل الهرمي المعدل لبولتزمان \eqref{eq:hier0}--\eqref{eq:hier_ell}، مع الحدود الإضافية المتناسبة مع $\beta\gamma$ والعمق البصري الفعال $\tau_{\mathrm{eff}}$. يُدرج عدم انحفاظ الطاقة-الزخم من $\dot{G}_{\mathrm{eff}}$ باتباع شكلية المرجع~\cite{Uzan2011}.
			\item \textbf{مقياس التخامد:} يُعاد حساب تكامل تخامد سيلك باستخدام سرعة الصوت واللزوجة المعدلتين.
			\item \textbf{الأنماط B الدورانية:} تحسب وحدة منفصلة طيف قدرة النمط $B$ من الدوران الشامل، مضافًا إلى مساهمات الموتر والعدسات القياسية.
		\end{enumerate}
		المعلمات الجديدة هي $\varepsilon_0$، $\eta$، وسعة الدوامية $A_v \propto \omega^2$. لتحليل MCMC، نثبت $\beta\gamma = 0.20$ (من الملحق P) ونعامل $\varepsilon_0$ و $\eta$ كمعلمتين حرتين. تُنوع المعلمات الكونية ($\Omega_b h^2$، $\Omega_c h^2$، $H_0$، $\tau_{\mathrm{reio}}$، $n_s$، $A_s$) كالمعتاد. نختبر أيضًا الارتباط بين المعلمات الجديدة و $n_s$، $A_s$ بفحص التوزيعات البعدية.
		
		\subsection{انحلالات المعلمات والنظاميات}
		
		تكشف التوزيعات البعدية عن ارتباط موجب خفيف بين $\varepsilon_0$ و $n_s$، وكذلك بين $\eta$ و $A_s$. هذا متوقع لأن تغيرًا في معدل التمدد المبكر أو مقياس التخامد يمكن تعويضه جزئيًا بتعديل الطيف البدئي. تبقى القيود المهمشة على معلمات $\Lambda$CDM القياسية دون تغيير جوهري، مما يظهر أن امتداد YUCT لا يقدم توترات قوية.
		
		التأثيرات المنهجية التي يمكن أن تحاكي فائض $\ell$ العالي تشمل تلوثًا متبقيًا من مصادر نقطية، ونمذجة غير دقيقة للحزمة، ومساهمات سنيياييف-زيلدوفيتش الحركية غير الخطية. من الضروري إعادة تحليل دقيق لبيانات Planck مع هذه التأثيرات لتأكيد حقيقة الإشارة. مع ذلك، يقدم نموذج YUCT تفسيرًا مدفوعًا فيزيائيًا (وإن كان غير قياسي) قابلًا للاختبار بالبيانات المستقبلية.
		
		\subsection{الدلالة الإحصائية وقيمة التحسن المتواضع}
		
		نلاحظ أن التحسن الإحصائي في $\chi^2$ متواضع ($\Delta\chi^2 = -7.6$ لمعلمتين إضافيتين، $\Delta\mathrm{BIC} \approx -2.2$). بمعزل عن غيره، لن يشكل هذا كشفًا عن فيزياء جديدة. لكن، لا ينبغي النظر إلى هذا التحليل ككشف مستقل. بدلاً من ذلك، يظهر أن التصحيحات \textit{الضرورية} التي تطلبها YUCT (الاعتماد على درجة الحرارة لـ $G_{\mathrm{eff}}$ ووجود اللزوجة التنسيقية) لا تفسد المواءمة الممتازة لـ $\Lambda$CDM مع البيانات. علاوة على ذلك، فهي تحسنها بشكل هامشي تمامًا في المنطقة حيث لوحظ فائض صغير. تكمن القوة الحقيقية لـ YUCT في المواءمة \textit{المتزامنة} لبيانات CMB هذه مع القيود المستقلة من انزياحات الأقزام البيضاء نحو الأحمر (الملحق~P)، ونظام الأرض-القمر (الملحق~P)، وانحرافات الرنين المعتمدة على الكتلة في الموجات الثقالية (الملحق~G). عندما تؤخذ كل مجموعات البيانات هذه بعين الاعتبار بشكل مشترك، تصبح الحجة لصالح الأس الكوني $\beta=0.67$ مقنعة.
		
		\subsection{تنبؤات للتجارب المستقبلية}
		
		سيكون CMB‑S4، بحساسيته المتوقعة البالغة $\sim 1\,\mu$K-arcmin من الضوضاء، قادرًا على كشف الأنماط $B$ الدورانية عند $>5\sigma$ إذا كانت معلمات YUCT قريبة من قيم المواءمة المثلى. سيوفر LiteBIRD فحصًا متقاطعًا حاسمًا عند مقاييس زاوية أكبر. علاوة على ذلك، يقع نمط اللا غاوسية المحدد ($f_{\mathrm{NL}}^{\mathrm{local}} \approx 1.1$) في متناول مسوحات البنى واسعة المقياس القادمة.
		
		\section{مناقشة وخاتمة}
		
		قدمنا أول معالجة كاملة لتباينات الخلفية الكونية الميكروية عند المقاييس الصغيرة في نظرية ياكوشيف الموحدة للتنسيق. النتائج الرئيسية هي:
		\begin{enumerate}
			\item \textbf{معادلات بولتزمان المعدلة} التي تدمج الجاذبية المعتمدة على درجة الحرارة واللزوجة التنسيقية، المشتقة من ازدواجات YUCT بين القطاعات ($\kappa_{01}$).
			\item \textbf{الاتصال بالحلقة الجبرية لـ YUCT}، مظهرين أن المعلمات الجديدة $\varepsilon_0$ و $\eta$ ليست اعتباطية بل مرتبطة بالثوابت الكونية $\beta=2/3$، $\kappa_c=1/3$، $S_{\mathrm{odd}}$، و $S_{\mathrm{even}}$. هذا يقلل العدد الفعال للمعلمات الحرة ويعزز الدافع الفيزيائي.
			\item \textbf{مواءمة محسنة مع بيانات Planck عند $\ell$ العالية}، مع $\Delta\chi^2 = -7.6$ لمعلمتين إضافيتين. يُفسر فائض القدرة عند $\ell \sim 2000$--$2500$ بشكل طبيعي بتخامد سيلك المخفض بسبب اللزوجة التنسيقية.
			\item \textbf{تنبؤ فريد بأنماط $B$ دورانية} ذات شكل $\ell^{-2}$ مميز، قابل للاختبار بتجارب CMB القادمة، ونمط محدد من اللا غاوسية البدئية.
		\end{enumerate}
		
		ناقشنا محدوديات التحليل الحالي: التحسن الإحصائي متواضع عند النظر إليه بمعزل، وتوجد انحلالات مع المعلمات القياسية. يمكن أن تساهم التأثيرات المنهجية في بيانات Planck أيضًا في بقايا $\ell$ العالي. مع ذلك، يقدم امتداد YUCT صورة فيزيائية متماسكة تربط الكون المبكر بفيزياء المختبر عبر الأس الكوني $\beta = 0.67$ والحلقة الجبرية. تكمن القوة الحقيقية لهذا الإطار في قدرته على تقديم تنبؤات متعددة ومترابطة عبر مجموعات بيانات متباينة (CMB، الأقزام البيضاء، الموجات الثقالية)، تنبع جميعها من مجموعة واحدة من الثوابت الكونية.
		
		يجب أن يتضمن العمل المستقبلي حسابًا أكثر تفصيلاً للزوجة التنسيقية من اللاغرانجيان الكامل ذي 19 بعدًا، بالإضافة إلى تحليلات MCMC مشتركة تجمع بين CMB والبنى واسعة المقياس والبيانات الفيزيائية الفلكية. يمكن أيضًا توسيع المنهجية المطورة هنا لتشمل رصدات البنى واسعة المقياس، موفرة اختبارًا شاملًا لنموذج التنسيق عبر جميع المقاييس الكونية.
		
		\section*{شكر وتقدير}
		
		يشكر المؤلف مطوري CAMB و MontePython على إتاحة أكوادهم علنًا.
		
		\section*{توفر البيانات}
		
		كود CAMB المعدل وسلاسل MCMC متاحة في مستودع YPSDC \url{https://github.com/Alexey-Yakushev-YUCT/YPSDC}.
		
		\begin{thebibliography}{99}
			\bibitem{Planck2018} Planck Collaboration, \emph{Astron. Astrophys.} \textbf{641}, A6 (2020).
			\bibitem{Ma1995} C.-P. Ma, E. Bertschinger, \emph{Astrophys. J.} \textbf{455}, 7 (1995).
			\bibitem{Silk1968} J. Silk, \emph{Astrophys. J.} \textbf{151}, 459 (1968).
			\bibitem{Dodelson2003} S. Dodelson, \emph{Modern Cosmology}, Academic Press (2003).
			\bibitem{Hu1995} W. Hu, N. Sugiyama, \emph{Astrophys. J.} \textbf{444}, 489 (1995).
			\bibitem{Cooray2005} A. Cooray, D. Baumann, \emph{Phys. Rev. D} \textbf{71}, 063002 (2005).
			\bibitem{Lewis2000} A. Lewis, A. Challinor, A. Lasenby, \emph{Astrophys. J.} \textbf{538}, 473 (2000).
			\bibitem{Audren2013} B. Audren et al., \emph{J. Cosmol. Astropart. Phys.} \textbf{1302}, 001 (2013).
			\bibitem{Uzan2011} J.-P. Uzan, \emph{Living Rev. Relativ.} \textbf{14}, 2 (2011).
			\bibitem{BICEP2021} BICEP/Keck Collaboration, \emph{Phys. Rev. Lett.} \textbf{127}, 151301 (2021).
			\bibitem{AppendixP} A.V. Yakushev, ``Appendix P: Temperature Dependence of Gravity,'' in \emph{YUCT}, Zenodo (2026).
			\bibitem{AppendixM} A.V. Yakushev, ``Appendix M: YUCT Cosmology — Inflation as a Coordination Phase Transition,'' in \emph{YUCT}, Zenodo (2026).
			\bibitem{AppendixΩ} A.V. Yakushev, ``Appendix $\Omega$: The Global Rotation of the Universe and Its New Minimum Age,'' in \emph{YUCT}, Zenodo (2026).
			\bibitem{AppendixA} A.V. Yakushev, ``Appendix A: Practical Guide to Using YUCT V35.0,'' in \emph{YUCT}, Zenodo (2026).
			\bibitem{AppendixY} A.V. Yakushev, ``Appendix Y: Mathematical Foundations of YUCT,'' in \emph{YUCT}, Zenodo (2026).
			\bibitem{AppendixLambda} A.V. Yakushev, ``Appendix $\Lambda$: Resolution of the Hierarchy and Cosmological Constant Problems within YUCT,'' in \emph{YUCT}, Zenodo (2026).
			\bibitem{AppendixFerra} A.V. Yakushev, ``Appendix Ferra1.2: Universal Coordination Constants for Ordered and Disordered Phases,'' in \emph{YUCT}, Zenodo (2026).
		\end{thebibliography}
		
	\end{multicols}
\end{document}